\begin{textsample}{POS Dim 3 – 2019\_c5 – Score 2.00 – 2019\_c5/enade\_2019\_eng\_mec\_9be5546a03.txt}  \label{ex:f3_pos_075}
CAPÍTULO 5 DISTRIBUIÇÃO DOS CONCEITOS 5.1 PANORAMA NACIONAL DA DISTRIBUIÇÃO DOS CONCEITOS O Conceito Enade é um indicador de qualidade que avalia os cursos por intermédio dos desempenhos dos estudantes a partir dos resultados do Enade.
Os valores possíveis vão de 1 ( pior situação ) a 5 ( melhor situação ).
Os cursos com apenas um ou sem nenhum concluinte participante não obtêm o Conceito Enade, ficando Sem Conceito ( SC ).
Na Tabela 5.1, são apresentadas a quantidade e a distribuição de cursos de Engenharia Mecânica participantes do Enade/2019, por faixa de conceito e Grande Região.
Enfatiza -se, mais uma vez, que a diferença entre os cursos tabulados neste capítulo e no capítulo 2 corresponde aos cursos que não foram avaliados, ou seja, àqueles, em princípio, sem estudantes concluintes que tivessem sido inscritos para a prova.
Observando -se os dados da Tabela 5.1, nota -se que, dos 385 cursos participantes, 163 ( 42,3 \% ) classificaram -se com conceito 2, o valor modal.
Esse foi também o conceito modal em todas as regiões : Norte ( 57,1 \% ), Nordeste ( 38,6 \% ), Sudeste ( 44,2 \% ), Sul ( 40,4 \% ) e Centro-Oeste ( 31,6 \% ).
O conceito 3 foi o segundo mais frequente em nível nacional ( 28,1 \%, correspondendo a 108 cursos ), e o conceito 4, o terceiro ( 13,0 \%, correspondendo a 50 cursos ).
Houve, ainda, 33 cursos ( 8,6 \% ) que receberam conceito 1, e 19 cursos ( 4,9 \% ) que receberam conceito 5.
Dos 385 cursos de Engenharia Mecânica, 12 ( 3,1 \% ) ficaram Sem Conceito ( SC ).
A região Norte participou com 14 cursos ou 3,6 \% do total nacional.
Desses, oito cursos receberam o conceito 2, o que equivale a 57,1 \% do total regional, além de ser o conceito modal na região, como já comentado.
A quatro cursos ( 28,6 \% ) atribuiu -se conceito 3, e a um curso ( 7,1 \% ), conceito 1.
Um curso ficou Sem Conceito ( SC ) e nenhum curso recebeu os conceitos 4 e 5 nessa região.
A região Nordeste participou com 57 cursos ou 14,8 \% do total nacional.
Nessa região, 22 cursos ( 38,6 \% em termos regionais ) obtiveram conceito 2, o conceito modal para a região, como já comentado.
Os conceitos 3 e 4 foram atribuídos, respectivamente, a 17 e a dez cursos ( 29,8 \% e 17,5 \%, respectivamente ).
Ainda nessa região, nenhum curso ficou com conceito 5.
Um curso ficou Sem Conceito ( SC ) e sete cursos ( 12,3 \% ) receberam conceito 1.
Dos 206 cursos participantes da região Sudeste ( 53,5 \% do total nacional ), o conceito 2 foi o conceito modal, como já anteriormente assinalado, tendo sido obtido por 91 cursos ( 44,2 \% ).
O conceito 1 foi atribuído a 21 cursos ( 10,2 \% ), e o conceito 3, a 47 cursos ( 22,8 \% ).
Receberam os conceitos 4 e 5, respectivamente, 28 cursos ( 13,6 \% ) e 12 cursos ( 5,8 \% ).
Nessa região, sete cursos ficaram Sem Conceito ( SC ).
Os 89 cursos da região Sul corresponderam a 23,1 \% do total nacional.
Como já indicado, houve predominância do conceito 2, atribuído a 40,4 \% dos cursos da região ( conceito modal ), ou seja, atribuído a 36 dos 89 cursos participantes na região Sul.
O conceito 3 foi atribuído a 35 cursos ( 39,3 \% ) e os conceitos 4 e 5, a nove ( 10,1 \% ) e cinco cursos ( 5,6 \% ), respectivamente.
Três cursos ficaram Sem Conceito ( SC ) na região Sul.
Já dos 19 cursos participantes na região Centro-Oeste ( 4,9 \% do total nacional ), seis ( 31,6 \% em termos regionais ) receberam conceito 2, o conceito modal, como já fora destacado.
Cinco cursos ( 26,3 \% ) receberam conceito 3, os conceitos 1 e 4 foram atribuídos a três cursos cada ( 15,8 \% cada ).
Dois cursos ( 10,5 \% ) receberam conceito 5 e nenhum ficou Sem Conceito ( SC ).
No Gráfico 5.1, é apresentada a distribuição cumulativa do Conceito Enade segundo Grande Região.
Em um gráfico de distribuição cumulativa, em cada valor de abscissa ( eixo x ), o valor apresentado no eixo das ordenadas ( eixo y ) é o valor acumulado da variável ( conceito Enade ).
Observando a linha vertical pontilhada, que \textbf{ajuda} a visualizar o percentual de conceitos acumulados até o valor 3, ou seja, conceitos 1, 2 e 3, pode -se notar que, por exemplo, a região Norte ( linha azul ) apresenta 100 \% dos cursos com conceito 3 ou abaixo deste valor.
A região Centro-Oeste ( linha \textbf{laranja} ), por outro lado, apresenta menos de 75 \% dos cursos na mesma situação : com conceito 3 ou abaixo.
Destaca -se que em algumas regiões o conceito 5 não foi alcançado por nenhum curso, como ocorre na região Norte, como já comentado, e na região Nordeste, onde todos os cursos têm conceito menor ou igual a 4, cuja poligonal cumulativa ( linha verde ) alcança 100 \% no conceito 4.
Em linhas gerais, em um gráfico de distribuição cumulativa, poligonais mais à esquerda representam Grandes Regiões com pior distribuição, e poligonais mais à direita, Grandes Regiões com melhores conceitos.
A região Norte ( linha azul ) apresenta os cursos com conceitos mais baixos, ao passo que a região Centro-Oeste ( linha \textbf{laranja} ) apresenta a distribuição com valores maiores.
Cursos em IES situadas nas regiões Sudeste ( linha vermelha ) e Sul ( linha rosa ) apresentam uma situação intermediária.
A região Norte ( linha azul ) apresenta os cursos com conceitos mais baixos, ao passo que a região Centro-Oeste ( linha \textbf{laranja} ) apresenta a distribuição com valores maiores.
Cursos em IES situadas nas regiões Sudeste ( linha vermelha ) e Sul ( linha rosa ) apresentam uma situação intermediária.
Em linhas gerais, em um gráfico de distribuição cumulativa, poligonais mais à esquerda representam Grandes Regiões com pior distribuição, e poligonais mais à direita, Grandes Regiões com melhores conceitos.
A região Norte ( linha azul ) apresenta os cursos com conceitos mais baixos, ao passo que a região Centro-Oeste ( linha \textbf{laranja} ) apresenta a distribuição com valores maiores.
Cursos em IES situadas nas regiões Sudeste ( linha vermelha ) e Sul ( linha rosa ) apresentam uma situação intermediária.
Grandes Regiões com melhores conceitos.
A região Norte ( linha azul ) apresenta os cursos com conceitos mais baixos, ao passo que a região Centro-Oeste ( linha \textbf{laranja} ) apresenta a distribuição com valores maiores.
Cursos em IES situadas nas regiões Sudeste ( linha vermelha ) e Sul ( linha rosa ) apresentam uma situação intermediária. 5.2 CONCEITOS POR CATEGORIA ADMINISTRATIVA, POR MODALIDADE DE ENSINO E GRANDE REGIÃO Na Tabela 5.2, é apresentada a distribuição dos cursos de Engenharia Mecânica participantes do Enade/2019, por Categoria Administrativa e Modalidade de ensino, de acordo com os conceitos por eles alcançados, segundo as Grandes Regiões brasileiras.
Dos 385 cursos participantes, 111 ( 28,8 \% ) eram ministrados em Instituições Públicas, e 274 ( 71,2 \% ), em Instituições Privadas.
Quanto à Modalidade de ensino, três cursos eram a Distância, e 382, Presenciais.
De acordo com as informações da Tabela 5.2, em termos nacionais, entre Instituições Públicas e Privadas, observa -se que, dos 19 cursos avaliados com conceito 5, 17 eram oferecidos em IES Públicas, e dois, em IES Privadas.
Dos 111 cursos participantes de IES Públicas, os conceitos 3 e 4 foram os valores modais, atribuídos a 42 cursos cada ( 37,8 \% cada ).
Dos demais cursos de Instituições Públicas participantes, dez cursos foram avaliados com conceito 2.
Nenhum curso ficou Sem Conceito ( SC ) ou recebeu conceito 1.
Como já comentado, 17 cursos ( 15,3 \% ), com conceito 5.
Na rede privada, o conceito modal foi 2, atribuído a 153 cursos, o correspondente a 55,8 \% dos 274 cursos da categoria.
Dos demais cursos participantes, 33 ( 12,0 \% ) receberam conceito 1, e 66 ( 24,1 \% ), conceito 3.
O conceito 4 foi atribuído a oito cursos ( 2,9 \% ), e o conceito 5, a dois cursos ( 0,7 \% ).
Nesta Categoria Administrativa, 12 cursos ( 4,4 \% ) ficaram Sem Conceito ( SC ).
Dos 382 cursos participantes da modalidade presencial, o conceito 2 foi o valor modal, atribuído a 162 cursos ( 42,4 \% ).
Dos demais cursos nesta modalidade de ensino, 33 cursos ( 8,6 \% ) receberam conceito 1, 108 cursos ( 28,3 \% ) receberam conceito 3, 50 cursos ( 13,1 \% ), conceito 4, e 19 cursos ( 5,0 \% ), conceito 5.
Nesta modalidade de ensino, dez cursos ( 2,6 \% ) ficaram Sem Conceito ( SC ).
Aos três únicos cursos a Distância informados, um curso recebeu conceito 2 e dois ficaram Sem Conceito ( SC ).
Na análise por região, observa -se que, na região Norte, as Instituições Privadas participaram com oito cursos ( 57,1 \% do total regional ), dos quais um ficou Sem Conceito ( SC ).
Nenhum curso recebeu os conceitos 3, 4 ou 5.
O conceito modal para as Instituições Privadas na região foi 2, atribuído a seis cursos, correspondendo a 75,0 \% dessa categoria na região.
Um curso foi avaliado com conceito 1.
As Instituições Públicas participaram com seis cursos na região Norte ( 42,9 \% do total regional ), aos quais foram atribuídos os conceitos 2 ( dois cursos ) e 3 ( quatro cursos ).
Nenhum curso ficou Sem Conceito ( SC ) ou com conceito 1, 4 e 5 nessa combinação de categoria e região. cursos ) e 3 ( quatro cursos ).
Nenhum curso ficou Sem Conceito ( SC ) ou com conceito 1, 4 e 5 nessa combinação de categoria e região.
Todos os cursos oferecidos na região Norte eram presenciais.
Na região Nordeste, a rede privada concentrou 33 dos 57 cursos participantes, o equivalente a 57,9 \% do total da região.
Dentre os cursos oferecidos por IES Privadas no Nordeste, 18 cursos foram avaliados com conceito 2, o conceito modal.
Sete cursos obtiveram conceitos 1 ou 3.
Nessa combinação de categoria e região, um curso ficou Sem Conceito ( SC ) e nenhum recebeu conceitos 4 e 5.
As Instituições Públicas da região Nordeste participaram com 24 cursos ( 42,1 \% do total da região ).
Desses, os conceitos 3 e 4 foram atribuídos a dez cursos cada.
Quatro cursos receberam conceito 2.
Nenhum curso ficou Sem Conceito ( SC ) ou com os demais conceitos nessa região.
Na região Nordeste, não houve cursos na Modalidade a Distância avaliados.
Na região Sudeste, a proporção de cursos da rede privada, 77,7 \%, foi mais elevada do que a registrada nas demais regiões brasileiras, correspondendo a 160 dos 206 cursos participantes.
Nessa categoria e região, o conceito modal foi 2, atribuído a 88 cursos, e sete cursos ficaram Sem Conceito ( SC ).
Os demais foram avaliados com conceito 1 ( 21 cursos ), conceito 3 ( 36 cursos ), conceito 4 ( sete cursos ) e conceito 5 ( um curso ).
Dos 46 cursos oferecidos por Instituições Públicas na região Sudeste, o conceito modal foi 4, atribuído a 21 cursos.
Nessa região, os conceitos 3 e 5 foram atribuídos a 11 cursos cada, três cursos receberam conceito 2.
Nenhum curso ficou Sem Conceito ( SC ) ou recebeu conceito 1.
Também na região Sudeste, a maioria dos cursos foi oferecida na modalidade Presencial : 203 cursos.
Para esses, a moda foi o conceito 2, com 90 cursos.
Os demais cursos presenciais da região foram alocados nos conceitos 1 ( 21 cursos ), 3 ( 47 cursos ), 4 ( 28 cursos ) e 5 ( doze cursos ).
Ficaram Sem Conceito ( SC ) cinco cursos.
Na região Sudeste, dois cursos oferecidos na modalidade Educação a Distância ficaram Sem Conceito ( SC ) e um curso recebeu conceito 2.
As Instituições Privadas concentraram 62 dos 89 cursos participantes da região Sul, 69,7 \% do total regional.
Desses, 36 cursos obtiveram conceito 2, o conceito modal.
Nesta combinação de Categoria Administrativa e Grande Região, 20 cursos receberam conceito 3.
Os conceitos 1, 4 e 5 foram atribuídos a um curso cada.
Três cursos ficaram Sem Conceito ( SC ).
As Instituições Públicas na região Sul participaram com 27 cursos ( 30,3 \% dos cursos da região ), aos quais foram atribuídos os conceitos 3 ( 15 cursos ), 4 ( oito cursos ) ou 5 ( quatro cursos ).
Nenhum curso da região Sul oferecido por IES Pública ficou Sem Conceito ( SC ) ou recebeu os conceitos 1 e 2.
Todos os cursos oferecidos na região Sul eram presenciais.
Na região Centro-Oeste, 11 dos 19 cursos participantes eram de Instituições Privadas ( 57,9 \% em termos regionais ).
Desses, cinco receberam conceito 2, o conceito modal.
Dos demais cursos, três receberam conceito 1, e outros três, conceito 3.
Nenhum curso ficou Sem Conceito ( SC ) ou recebeu conceitos 4 ou 5.
Quanto aos oito cursos oferecidos por Instituições Públicas na região Centro-Oeste ( 42,1 \% do total regional ), três cursos apresentaram conceito 4, conceito modal.
Os conceitos 3 e 5 foram atribuídos a dois cursos cada, e um curso recebeu conceito 2.
Nenhum curso ficou Sem Conceito ( SC ) ou recebeu conceito 1.
Todos os cursos oferecidos na região Centro-Oeste eram presenciais.
No Gráfico 5.2, apresenta -se a distribuição do Conceito Enade, segundo a Categoria Administrativa da IES.
Os cursos em IES Públicas ( linha azul ) apresentam uma poligonal mais à direita do que os em IES Privadas ( linha verde ) e, consequentemente, uma distribuição de Conceitos Enade com valores maiores. 5.3 CONCEITOS POR ORGANIZAÇÃO ACADÊMICA E GRANDE REGIÃO Na Tabela 5.3, encontra -se a distribuição dos conceitos atribuídos aos cursos participantes do Enade/2019, na Área de Engenharia Mecânica, por Organização Acadêmica, segundo as Grandes Regiões brasileiras.
Dos 385 cursos de Engenharia Mecânica participantes, 174 eram oferecidos em Universidades, 89 em Centros Universitários, 96 em Faculdades e 26 em CEFET/IF.
Esta distribuição corresponde a, respectivamente, 45,2 \%, 23,1 \%, 24,9 \% e 6,8 \% dos cursos.
De acordo com os dados apresentados, de todos os 19 cursos avaliados com conceito 5, nove eram vinculados a Universidades.
Os cursos vinculados a esse tipo de Organização Acadêmica tiveram o conceito 3 como conceito modal, atribuído a 69 cursos ( 39,7 \% ).
Os demais cursos vinculados a Universidades receberam os conceitos 1 ( cinco cursos ), 2 ( 53 cursos ), 4 ( 35 cursos ) e conceito 5 ( nove cursos, como já mencionado ).
Ficaram Sem Conceito ( SC ), três cursos.
Dos cursos em Centros Universitários, o conceito modal foi 2, atribuído a 55 cursos ( 61,8 \% ).
Dois cursos vinculados a esse tipo de Organização Acadêmica ficaram Sem Conceito ( SC ), e os outros cursos receberam os conceitos 1 ( oito cursos ), 3 ( 19 cursos ), 4 ( dois cursos ) e 5 ( três cursos ).
Dos 96 cursos mantidos por Faculdades, 52 ( 54,2 \% ) receberam conceito 2, o conceito modal.
O segundo conceito mais frequentemente atribuído aos cursos oferecidos por essa categoria de IES foi o conceito 1, recebido por 20 cursos.
Na sequência, nove cursos obtiveram conceito 3, os conceitos 4 e 5 foram atribuídos a quatro cursos cada.
Dos cursos oferecidos em Faculdades, sete ficaram Sem Conceito ( SC ).
Dos 26 cursos mantidos por CEFET/IF, 11 ( 42,3 \% ) receberam conceito 3, o conceito modal.
O segundo conceito mais frequentemente atribuído aos cursos oferecidos por essa categoria de IES foi o conceito 4, recebido por nove cursos.
Os conceitos 2 e 5 foram atribuídos a três cursos cada.
Nenhum curso ficou Sem Conceito ( SC ), e nenhum curso recebeu conceito 1.
Considerando -se, separadamente, as regiões brasileiras, verifica -se que, na região Norte, as Universidades participaram com sete dos 14 cursos avaliados, aos quais foram atribuídos conceito 2 ( três cursos ) e 3 ( quatro cursos ).
Nenhum curso ficou Sem Conceito ( SC ) ou recebeu os demais conceitos nesta combinação de categoria e região.
Os Centros Universitários da região Norte foram representados por três cursos, os quais obtiveram conceitos 1 e 2.
Um curso ficou Sem Conceito ( SC ).
Nenhum curso recebeu os demais conceitos nesta combinação de categoria e região.
As Faculdades participaram com três cursos na região Norte, os quais receberam o conceito 2.
Já os CEFET/IF participaram com um curso na região Norte, ao qual foi atribuído o conceito 2 Na região Nordeste, as Universidades participaram com 26 dos 57 cursos da Área de Engenharia Mecânica oferecidos.
Dez cursos receberam conceito 3, nove cursos receberam conceito 4 e sete cursos receberam conceito 2.
Nenhum curso ficou Sem Conceito ( SC ) ou recebeu os demais conceitos.
Os Centros Universitários contaram com 12 cursos participantes na região Nordeste, dez dos quais receberam o conceito modal 2.
Para os demais cursos, foi atribuído o conceito 3.
Nenhum curso oferecido por Centros Universitários no Nordeste ficou Sem Conceito ( SC ) ou recebeu os conceitos 1, 4 e 5.
As Faculdades foram representadas por 14 cursos na região Nordeste, sete dos quais receberam conceito 1, o conceito modal.
Quatro cursos receberam conceito 2 e dois receberam conceito 3.
Um curso ficou Sem Conceito ( SC ) e nenhum recebeu os conceitos 4 ou 5.
Os CEFET/IF participaram com cinco cursos na região Nordeste, avaliados com conceitos 3 ( três cursos ), 2 e 4 ( um curso cada ).
Nenhum curso ficou Sem Conceito ( SC ), ou recebeu os conceitos 1 ou 5.
Na região Sudeste, as Universidades concentraram 84 dos 206 cursos de Engenharia Mecânica da região.
Dos cursos oferecidos em Universidades, nessa região, os conceitos modais foram 2 e 3, atribuídos a 29 cursos cada.
Os demais cursos receberam os conceitos 1 ( quatro cursos ), 4 ( 16 cursos ) e 5 ( três cursos ).
Três cursos ficaram Sem Conceito ( SC ).
Os Centros Universitários participaram com 51 cursos na região Sudeste, dos quais 28 obtiveram o conceito modal 2.
O conceito 1 foi atribuído a sete cursos, o conceito 3, a dez cursos, conceito 4, a dois cursos, e o conceito 5, a três cursos.
Um curso ficou Sem Conceito ( SC ) nessa combinação de Organização Acadêmica e região.
Já as Faculdades foram representadas por 59 cursos na região Sudeste, dos quais três ficaram Sem Conceito ( SC ).
O conceito modal 2 foi atribuído a 33 cursos.
Aos demais cursos foram atribuídos os conceitos 1 ( dez cursos ), 3 ( cinco cursos ), 4 ( quatro cursos ) e 5 ( quatro cursos ).
Os CEFET/IF participaram com 12 cursos, seis deles receberam conceito 4, três cursos receberam conceito 3, dois cursos receberam conceito 5 e um curso recebeu conceito 2.
Nenhum curso ficou Sem Conceito ( SC ) ou foi avaliado com conceito 1.
Dos 89 cursos da região Sul, 47 eram mantidos por Universidades, dos quais o conceito modal foi 3, atribuído a 23 cursos.
O segundo conceito mais frequente foi 2 ( 12 cursos ), seguido do conceito 4 ( sete cursos ).
Quatro cursos receberam conceito 5 e um curso recebeu conceito 1.
Nenhum curso ficou Sem Conceito ( SC ).
Dos 19 cursos participantes de Centros Universitários da região Sul, 12 receberam conceito 2, e os outros sete receberam conceito 3.
Nenhum curso ficou Sem Conceito ( SC ) ou recebeu os demais conceitos.
Quanto aos 16 cursos vinculados a Faculdades na região Sul, 12 cursos receberam o conceito modal 2, e um curso recebeu conceito 3.
Nenhum curso recebeu os demais conceitos, e três cursos ficaram Sem Conceito ( SC ).
Já os CEFET/IF participaram com sete cursos na região Centro-Oeste, aos quais foram atribuídos os conceitos 3 ( quatro cursos ), 4 ( dois cursos ) e 5 ( um curso ).
Na região Centro-Oeste, dez dos 19 cursos oferecidos eram mantidos por Universidades.
Para este tipo de Organização Acadêmica, os conceitos 3 e 4 foram atribuídos a três cursos cada, já os conceitos 2 e 5 foram atribuídos a dois cursos cada.
Nenhum curso ficou Sem Conceito ( SC ) ou recebeu conceito 1.
Os Centros Universitários da região Centro-Oeste contaram com quatro cursos, todos os quais avaliados com conceito 2.
Dos quatro cursos oferecidos em Faculdades na região Centro-Oeste, três receberam o conceito modal 1, e um obteve conceito 3.
Já os CEFET/IF participaram com um curso na região Centro-Oeste, ao qual foi atribuído o conceito 3.
No Gráfico 5.3, apresenta -se a distribuição do Conceito Enade dos cursos de Engenharia Mecânica, segundo a Organização Acadêmica da IES.
Os cursos em CEFET/IF ( linha \textbf{laranja} ) apresentam uma poligonal mais à direita do que as demais, denotando uma melhor distribuição de Conceitos Enade.
No outro extremo, com distribuição muito semelhante entre si, encontram -se os cursos mantidos por Centros Universitários ( linha verde ) e Faculdades ( linha vermelha ), apresentando as respectivas poligonais mais à esquerda e, consequentemente, distribuição de Conceitos Enade com menores valores.

% matched lemmas: ajudar, laranja
\end{textsample}
